\chapter{The Singular Value Decomposition}

The spectral theorem diagonalises symmetric matrices. What about everything else: rectangular matrices, non-symmetric matrices, matrices with missing eigenvectors? The \emph{singular value decomposition} (SVD) says that every matrix, no matter its shape or rank, acts as a rotation, then a stretch along coordinate axes, then another rotation. It exists for every matrix, it is numerically stable, and it reveals rank, the four fundamental subspaces, the best low-rank approximation, the pseudo-inverse and the condition number all at once. It is the single most useful matrix factorisation, and the foundation of principal component analysis, latent semantic analysis, recommender systems, image compression and low-rank adapters in machine learning.

\section{The theorem}

\begin{theorem}[Singular value decomposition]\label{thm:svd}
Every matrix $A\in\R^{m\times n}$ of rank $r$ can be written as
\[ A=U\Sigma V\T, \]
where $U\in\R^{m\times m}$ and $V\in\R^{n\times n}$ are orthogonal and $\Sigma\in\R^{m\times n}$ has $\Sigma_{ii}=\sigma_i$ for $i=1,\dots,\min(m,n)$ and zeros elsewhere, with
\[ \sigma_1\ge\sigma_2\ge\dots\ge\sigma_r>0=\sigma_{r+1}=\dots. \]
The numbers $\sigma_i$ are the \emph{singular values} of $A$ (uniquely determined); the columns $\mathbf{u}_i$ of $U$ are the \emph{left singular vectors} and the columns $\mathbf{v}_i$ of $V$ the \emph{right singular vectors}. Equivalently
\[ A=\sum_{i=1}^r\sigma_i\,\mathbf{u}_i\mathbf{v}_i\T,\qquad A\mathbf{v}_i=\sigma_i\mathbf{u}_i,\qquad A\T\mathbf{u}_i=\sigma_i\mathbf{v}_i. \]
\end{theorem}

\begin{proof}
The matrix $A\T A\in\R^{n\times n}$ is symmetric positive semidefinite, so by the spectral theorem it has an orthonormal eigenbasis $\mathbf{v}_1,\dots,\mathbf{v}_n$ with eigenvalues $\lambda_1\ge\dots\ge\lambda_n\ge0$. Since $\rank A\T A=\rank A=r$ (Proposition~\ref{prop:rankineq}), exactly $r$ eigenvalues are positive. Put $\sigma_i=\sqrt{\lambda_i}$. For $i\le r$ define $\mathbf{u}_i=A\mathbf{v}_i/\sigma_i$. These are orthonormal:
\[ \ip{\mathbf{u}_i}{\mathbf{u}_j}=\frac{\mathbf{v}_i\T A\T A\mathbf{v}_j}{\sigma_i\sigma_j}=\frac{\lambda_j\mathbf{v}_i\T\mathbf{v}_j}{\sigma_i\sigma_j}=\delta_{ij}. \]
For $i>r$ we have $\norm{A\mathbf{v}_i}^2=\lambda_i=0$, so $A\mathbf{v}_i=\mathbf{0}$. Extend $\mathbf{u}_1,\dots,\mathbf{u}_r$ to an orthonormal basis $\mathbf{u}_1,\dots,\mathbf{u}_m$ of $\R^m$. Then $AV=U\Sigma$ (check column by column), hence $A=U\Sigma V\T$. Also $A\T\mathbf{u}_i=A\T A\mathbf{v}_i/\sigma_i=\sigma_i\mathbf{v}_i$. Uniqueness of the $\sigma_i$: $A\T A=V\Sigma\T\Sigma V\T$ shows $\sigma_i^2$ are the eigenvalues of $A\T A$.
\end{proof}

\begin{definition}[Compact (thin) SVD]
Keeping only the nonzero singular values gives $A=U_r\Sigma_rV_r\T$ with $U_r\in\R^{m\times r}$ and $V_r\in\R^{n\times r}$ with orthonormal columns and $\Sigma_r=\diag(\sigma_1,\dots,\sigma_r)$ square and invertible. This is the form stored in practice.
\end{definition}

\begin{example}[Computing an SVD]\label{ex:svd}
Let $A=\begin{pmatrix}3&0\\4&5\end{pmatrix}$. Then $A\T A=\begin{pmatrix}25&20\\20&25\end{pmatrix}$ has eigenvalues $45$ and $5$ with eigenvectors $\mathbf{v}_1=\frac1{\sqrt2}(1,1)$ and $\mathbf{v}_2=\frac1{\sqrt2}(-1,1)$. So $\sigma_1=\sqrt{45}=3\sqrt5$, $\sigma_2=\sqrt5$. Next $\mathbf{u}_1=A\mathbf{v}_1/\sigma_1=\frac1{\sqrt2}(3,9)/(3\sqrt5)=\frac1{\sqrt{10}}(1,3)$ and $\mathbf{u}_2=A\mathbf{v}_2/\sigma_2=\frac1{\sqrt2}(-3,1)/\sqrt5=\frac1{\sqrt{10}}(-3,1)$. A check: $\sigma_1\sigma_2=15=\abs{\det A}$, and $\sigma_1^2+\sigma_2^2=50=\norm{A}_F^2$.
\end{example}

\section{Geometry: spheres become ellipsoids}

Reading $A=U\Sigma V\T$ right to left: $V\T$ rotates (or reflects) the input; $\Sigma$ stretches the $i$-th coordinate axis by $\sigma_i$ (and maps $\R^n\to\R^m$, dropping or padding with zeros); $U$ rotates the result into place.

\begin{figure}[h]
\centering
\begin{tikzpicture}[scale=1.0]
  \begin{scope}
    \draw[axis] (-1.6,0)--(1.6,0); \draw[axis] (0,-1.6)--(0,1.6);
    \draw[thick,accent2] (0,0) circle (1.1);
    \begin{scope}[rotate=45]\draw[vec,accent] (0,0)--(1.1,0) node[above right,lab]{$\mathbf{v}_1$}; \draw[vec,accent3] (0,0)--(0,1.1) node[above left,lab]{$\mathbf{v}_2$};\end{scope}
    \node[lab] at (0,-2.0){unit circle};
  \end{scope}
  \draw[->,thick,mid] (2.0,0)--(3.4,0) node[midway,above,lab]{$A$};
  \begin{scope}[xshift=6cm]
    \draw[axis] (-2.8,0)--(2.8,0); \draw[axis] (0,-1.6)--(0,1.6);
    \begin{scope}[rotate=20]
      \draw[thick,accent2] (0,0) ellipse (2.4cm and 0.6cm);
      \draw[vec,accent] (0,0)--(2.4,0) node[above right,lab]{$\sigma_1\mathbf{u}_1$};
      \draw[vec,accent3] (0,0)--(0,0.6) node[above left,lab]{$\sigma_2\mathbf{u}_2$};
    \end{scope}
    \node[lab] at (0,-2.0){image: an ellipse};
  \end{scope}
\end{tikzpicture}
\caption{A matrix maps the unit circle to an ellipse. The semi-axes of the ellipse are $\sigma_i\mathbf{u}_i$, the images of the right singular vectors $\mathbf{v}_i$.}
\end{figure}

\begin{theorem}[SVD and the four subspaces]\label{thm:svd4}
With the notation of Theorem~\ref{thm:svd}:
\begin{align*}
\Col A&=\Span\{\mathbf{u}_1,\dots,\mathbf{u}_r\}, & \Null A\T&=\Span\{\mathbf{u}_{r+1},\dots,\mathbf{u}_m\},\\
\Row A&=\Span\{\mathbf{v}_1,\dots,\mathbf{v}_r\}, & \Null A&=\Span\{\mathbf{v}_{r+1},\dots,\mathbf{v}_n\}.
\end{align*}
So the SVD provides orthonormal bases of all four fundamental subspaces simultaneously, and $\rank A$ is the number of nonzero singular values.
\end{theorem}

\begin{proof}
$A\mathbf{x}=\sum_{i\le r}\sigma_i(\mathbf{v}_i\T\mathbf{x})\mathbf{u}_i$ shows $\Col A\subseteq\Span\{\mathbf{u}_1,\dots,\mathbf{u}_r\}$, and $A\mathbf{v}_i=\sigma_i\mathbf{u}_i$ gives the reverse inclusion. $A\mathbf{x}=\mathbf{0}$ iff $\mathbf{v}_i\T\mathbf{x}=0$ for $i\le r$, giving $\Null A$. The statements for $A\T=V\Sigma\T U\T$ are identical with the roles of $U,V$ exchanged.
\end{proof}

\section{Norms and the condition number}

\begin{definition}[Matrix norms]
The \emph{spectral norm} is $\norm{A}_2=\max_{\norm{\mathbf{x}}=1}\norm{A\mathbf{x}}$, the largest stretching factor. The \emph{Frobenius norm} is $\norm{A}_F=\sqrt{\sum a_{ij}^2}$. The \emph{nuclear norm} is $\norm{A}_*=\sum\sigma_i$.
\end{definition}

\begin{theorem}[Norms via singular values]\label{thm:svdnorms}
$\norm{A}_2=\sigma_1$, $\norm{A}_F=\sqrt{\sigma_1^2+\dots+\sigma_r^2}$ and $\norm{A}_*=\sigma_1+\dots+\sigma_r$. All three are invariant under multiplication by orthogonal matrices.
\end{theorem}

\begin{proof}
$\norm{A\mathbf{x}}^2=\mathbf{x}\T A\T A\mathbf{x}$ is a Rayleigh quotient of $A\T A$ on unit vectors, with maximum $\lambda_1=\sigma_1^2$ (Theorem~\ref{thm:rayleigh}). $\norm{A}_F^2=\tr(A\T A)=\sum\lambda_i=\sum\sigma_i^2$ by Corollary~\ref{cor:trdet}. Invariance follows because $A$ and $QAQ'$ have the same singular values.
\end{proof}

\begin{definition}[Condition number]
For invertible $A$ the (spectral) \emph{condition number} is $\kappa(A)=\norm{A}_2\norm{A^{-1}}_2=\sigma_1/\sigma_n$.
\end{definition}

\begin{theorem}[Sensitivity of linear systems]\label{thm:cond}
If $A\mathbf{x}=\mathbf{b}$ and $A(\mathbf{x}+\delta\mathbf{x})=\mathbf{b}+\delta\mathbf{b}$, then
\[ \frac{\norm{\delta\mathbf{x}}}{\norm{\mathbf{x}}}\le\kappa(A)\frac{\norm{\delta\mathbf{b}}}{\norm{\mathbf{b}}}. \]
\end{theorem}

\begin{proof}
$\delta\mathbf{x}=A^{-1}\delta\mathbf{b}$ so $\norm{\delta\mathbf{x}}\le\norm{A^{-1}}\norm{\delta\mathbf{b}}$, and $\norm{\mathbf{b}}=\norm{A\mathbf{x}}\le\norm{A}\norm{\mathbf{x}}$. Divide.
\end{proof}

A condition number near $1$ means well-conditioned; $\kappa=10^{k}$ means one may lose up to $k$ digits of accuracy. Orthogonal matrices have $\kappa=1$. Chapter~18 develops these ideas.

\section{Low-rank approximation: Eckart--Young}

The compact SVD writes $A$ as a sum of $r$ rank-one pieces ordered by importance $\sigma_i$. Truncating after $k$ terms gives $A_k=\sum_{i\le k}\sigma_i\mathbf{u}_i\mathbf{v}_i\T$.

\begin{theorem}[Eckart--Young--Mirsky]\label{thm:eckart}
For every matrix $B$ of rank at most $k$,
\[ \norm{A-B}_2\ge\sigma_{k+1}=\norm{A-A_k}_2,\qquad\norm{A-B}_F\ge\sqrt{\sigma_{k+1}^2+\dots+\sigma_r^2}=\norm{A-A_k}_F. \]
Thus $A_k$ is the best rank-$k$ approximation of $A$ in both norms.
\end{theorem}

\begin{proof}[Proof for the spectral norm]
$\norm{A-A_k}_2=\norm{\sum_{i>k}\sigma_i\mathbf{u}_i\mathbf{v}_i\T}_2=\sigma_{k+1}$. Now let $\rank B\le k$; then $\dim\Null B\ge n-k$. The subspace $\Span\{\mathbf{v}_1,\dots,\mathbf{v}_{k+1}\}$ has dimension $k+1$, so it meets $\Null B$ nontrivially (dimensions $k+1$ and $n-k$ add to $n+1>n$, so the intersection has dimension $\ge1$). Take a unit vector $\mathbf{x}=\sum_{i\le k+1}c_i\mathbf{v}_i$ in the intersection. Then $B\mathbf{x}=\mathbf{0}$, so
\[ \norm{(A-B)\mathbf{x}}^2=\norm{A\mathbf{x}}^2=\sum_{i\le k+1}\sigma_i^2c_i^2\ge\sigma_{k+1}^2\sum c_i^2=\sigma_{k+1}^2. \]
Hence $\norm{A-B}_2\ge\sigma_{k+1}$. The Frobenius case is proved similarly using Weyl's inequality for singular values.
\end{proof}

\begin{intuition}
The first few singular values often capture most of the ``energy'' $\norm{A}_F^2=\sum\sigma_i^2$. A grayscale image is a matrix; keeping $k$ singular triples stores $k(m+n+1)$ numbers instead of $mn$. For a $1000\times1000$ image, $k=50$ stores about $10\%$ of the data yet often looks nearly identical. Latent semantic analysis, recommender systems and denoising work the same way: small singular values are mostly noise.
\end{intuition}

\begin{interview}{choosing the rank}
A common recipe: keep the smallest $k$ with $\sum_{i\le k}\sigma_i^2/\sum_i\sigma_i^2\ge0.9$ (``90\% of variance''), or look for an elbow in the plot of $\sigma_i$ against $i$ (the \emph{scree plot}). Eckart--Young guarantees that no other rank-$k$ matrix is closer, so there is no better compression of that rank to look for.
\end{interview}

\section{The pseudo-inverse}

When $A$ is not square or not invertible, $A^{-1}$ does not exist; the SVD provides the best substitute.

\begin{definition}[Moore--Penrose pseudo-inverse]
$A^+=V\Sigma^+U\T$, where $\Sigma^+\in\R^{n\times m}$ has $\sigma_i^{-1}$ in the first $r$ diagonal positions and zeros elsewhere. Equivalently $A^+=\sum_{i\le r}\sigma_i^{-1}\mathbf{v}_i\mathbf{u}_i\T$.
\end{definition}

\begin{theorem}[Least squares via $A^+$]\label{thm:pinv}
For every $A$ and $\mathbf{b}$, the vector $\hat{\mathbf{x}}=A^+\mathbf{b}$ is the least squares solution of smallest norm: it minimises $\norm{A\mathbf{x}-\mathbf{b}}$, and among all minimisers it has the smallest $\norm{\mathbf{x}}$. Moreover: $A^+=(A\T A)^{-1}A\T$ if $A$ has independent columns; $A^+=A\T(AA\T)^{-1}$ if $A$ has independent rows; $A^+=A^{-1}$ if $A$ is invertible; $AA^+$ is the orthogonal projection onto $\Col A$ and $A^+A$ onto $\Row A$.
\end{theorem}

\begin{proof}
Substitute the SVD: $\norm{A\mathbf{x}-\mathbf{b}}^2=\norm{\Sigma\mathbf{y}-\mathbf{c}}^2$ with $\mathbf{y}=V\T\mathbf{x}$, $\mathbf{c}=U\T\mathbf{b}$ (orthogonal matrices preserve norms). This equals $\sum_{i\le r}(\sigma_iy_i-c_i)^2+\sum_{i>r}c_i^2$, minimised exactly when $y_i=c_i/\sigma_i$ for $i\le r$, with $y_i$ free for $i>r$. Since $\norm{\mathbf{x}}=\norm{\mathbf{y}}$, the smallest-norm choice sets $y_i=0$ for $i>r$, i.e.\ $\mathbf{y}=\Sigma^+\mathbf{c}$ and $\mathbf{x}=V\Sigma^+U\T\mathbf{b}=A^+\mathbf{b}$. The special formulas are checked by substituting $A=U_r\Sigma_rV_r\T$.
\end{proof}

\begin{example}
For the rank-one matrix $A=\begin{pmatrix}1&1\\1&1\end{pmatrix}$: $\sigma_1=2$, $\mathbf{v}_1=\mathbf{u}_1=\frac1{\sqrt2}(1,1)$, so $A^+=\frac12\mathbf{v}_1\mathbf{u}_1\T=\frac14\begin{pmatrix}1&1\\1&1\end{pmatrix}$. The system $A\mathbf{x}=(2,2)\T$ has solutions on the line $x_1+x_2=2$; $A^+(2,2)\T=(1,1)\T$ is the one of smallest norm.
\end{example}

\section{Ridge regression, PCA and the polar decomposition}

\begin{proposition}[Ridge in the SVD basis]\label{prop:ridgesvd}
The ridge solution $(A\T A+\lambda I)^{-1}A\T\mathbf{b}$ equals $\sum_{i\le r}\dfrac{\sigma_i}{\sigma_i^2+\lambda}(\mathbf{u}_i\T\mathbf{b})\mathbf{v}_i$. Compared with the unregularised solution $\sum\frac1{\sigma_i}(\mathbf{u}_i\T\mathbf{b})\mathbf{v}_i$, each component is multiplied by the \emph{filter factor} $\sigma_i^2/(\sigma_i^2+\lambda)\in(0,1)$, which is close to $1$ for large $\sigma_i$ and close to $0$ for small $\sigma_i$.
\end{proposition}

\begin{proof}
Substituting $A=U\Sigma V\T$: $A\T A+\lambda I=V(\Sigma\T\Sigma+\lambda I)V\T$ and $A\T\mathbf{b}=V\Sigma\T U\T\mathbf{b}$; invert the diagonal middle factor.
\end{proof}

\begin{theorem}[PCA via SVD]\label{thm:pca}
Let $X\in\R^{m\times n}$ have centred columns (each feature has mean $0$) and $X=U\Sigma V\T$. The principal components (directions of maximal variance) are the columns $\mathbf{v}_i$ of $V$; the variance along $\mathbf{v}_i$ is $\sigma_i^2/(m-1)$; the scores (coordinates of the data in the principal basis) are $XV=U\Sigma$.
\end{theorem}

\begin{proof}
The sample covariance matrix is $C=X\T X/(m-1)=V\Sigma\T\Sigma V\T/(m-1)$, so its eigenvectors are the $\mathbf{v}_i$ with eigenvalues $\sigma_i^2/(m-1)$; the rest is Theorem~\ref{thm:rayleigh}. Projecting onto the first $k$ components is exactly the rank-$k$ truncation $X_k$.
\end{proof}

\begin{theorem}[Polar decomposition]\label{thm:polar}
Every square matrix factors as $A=WP$ with $W$ orthogonal and $P=(A\T A)^{1/2}\succeq0$ (a ``rotation times a stretch''), with $P$ unique and $W$ unique if $A$ is invertible.
\end{theorem}

\begin{proof}
From $A=U\Sigma V\T$ set $W=UV\T$ and $P=V\Sigma V\T$; then $WP=UV\T V\Sigma V\T=U\Sigma V\T=A$, and $P^2=V\Sigma^2V\T=A\T A$.
\end{proof}

\begin{nuance}{signs and uniqueness}
The singular values are unique, but singular vectors are not. $\mathbf{v}_i$ and $\mathbf{u}_i$ can be flipped together ($\mathbf{u}_i\to-\mathbf{u}_i$, $\mathbf{v}_i\to-\mathbf{v}_i$) without changing $A$, and if a singular value repeats, any orthonormal basis of the corresponding singular subspace can be used. Software may return different signs on different platforms; never compare singular vectors directly, compare the subspaces or the products $\mathbf{u}_i\mathbf{v}_i\T$.
\end{nuance}

\begin{summarybox}
\begin{itemize}
  \item Every matrix has an SVD $A=U\Sigma V\T$ with orthogonal $U,V$ and nonnegative diagonal $\Sigma$; singular values are $\sigma_i=\sqrt{\lambda_i(A\T A)}$ and $\rank A$ counts the nonzero ones.
  \item $A$ maps the unit sphere to an ellipsoid with semi-axes $\sigma_i\mathbf{u}_i$; the SVD gives orthonormal bases of all four fundamental subspaces.
  \item $\norm{A}_2=\sigma_1$, $\norm{A}_F^2=\sum\sigma_i^2$, $\kappa(A)=\sigma_1/\sigma_n$.
  \item Eckart--Young: truncating the SVD at rank $k$ gives the best rank-$k$ approximation in the spectral and Frobenius norms.
  \item The pseudo-inverse $A^+=V\Sigma^+U\T$ gives the minimum-norm least squares solution.
  \item PCA is the SVD of centred data; ridge shrinks each singular direction by $\sigma_i^2/(\sigma_i^2+\lambda)$.
\end{itemize}
\end{summarybox}

\section{Exercises}

\begin{exercise}
Compute the SVD of $\begin{pmatrix}1&1\\0&0\\1&-1\end{pmatrix}$ and identify its rank, nullspace and best rank-one approximation.
\end{exercise}

\begin{exercise}
Show that the singular values of a symmetric matrix are the absolute values of its eigenvalues, and those of an orthogonal matrix are all $1$.
\end{exercise}

\begin{exercise}
Prove that $\sigma_1\ge\abs{\lambda}$ for every eigenvalue $\lambda$ of a square matrix $A$ and that $\abs{\det A}=\prod\sigma_i$.
\end{exercise}

\begin{exercise}
Find the pseudo-inverse of $\begin{pmatrix}2&0\\0&0\\0&3\end{pmatrix}$ and of $(1\ 2\ 2)$, and solve $(1\ 2\ 2)\mathbf{x}=3$ in the minimum-norm sense.
\end{exercise}

\begin{exercise}
Let $A$ be $100\times 50$ with singular values $10,5,1,0.1,\dots,0.001$. What are $\norm{A}_2$, $\kappa$ (if invertible), and the error $\norm{A-A_2}_2$?
\end{exercise}

\begin{exercise}
Prove that if $\rank B\le k$ then $\norm{A-B}_F^2\ge\sum_{i>k}\sigma_i^2$, using Weyl's inequality $\sigma_{i+j-1}(X+Y)\le\sigma_i(X)+\sigma_j(Y)$.
\end{exercise}

\begin{exercise}
Show that $\norm{A}_2\le\norm{A}_F\le\sqrt{\rank A}\,\norm{A}_2$ and give matrices achieving each equality.
\end{exercise}

\begin{exercise}
Explain how to obtain the SVD of $A$ from the eigendecomposition of the symmetric matrix $\begin{pmatrix}0&A\\A\T&0\end{pmatrix}$, whose eigenvalues are $\pm\sigma_i$ (and zeros).
\end{exercise}
